\chapter{群、群作用与表示论}
\label{chap:group-theory}

群论在数学物理中的作用不是给对称性对象贴上名称，而是把“可复合、可逆的变换”变成可计算的代数结构。真正进入量子力学、谱理论和几何以后，仅知道群的乘法还不够：还需要知道群如何作用在集合、向量空间和算子空间上，以及一个表示如何分解成不可约成分。因此本章从群与商结构开始，依次建立群作用、有限群表示、character 与投影算符，最后进入 Lie 群、Lie 代数、伴随表示和 Casimir 元。循环群、二面体群、\(SU(2)\) 等熟悉对象只在一般理论建立以后作为特例出现。

\section{群、商结构与群作用}
\label{sec:group-basics}

接下来考察群、同态与第一同构定理。

设 \(G\) 是集合，二元运算写成乘法。若乘法满足结合律，存在单位元 \(e\)，并且每个 \(g\in G\) 都存在逆元 \(g^{-1}\)，则称 \(G\) 为群。若还满足 \(gh=hg\)，则称为 Abel 群。有限群的元素个数记为 \(|G|\)。由一个元素 \(g\) 生成的循环子群记为
\[
\langle g\rangle=\{g^n:n\in\mathbb Z\}.
\]
若存在最小正整数 \(m\) 使 \(g^m=e\)，则称 \(m\) 为 \(g\) 的阶，记为 \(\operatorname{ord}(g)=m\)。

群之间保持乘法的映射
\[
\varphi:G\longrightarrow H,
\qquad
\varphi(g_1g_2)=\varphi(g_1)\varphi(g_2)
\]
称为群同态。其核与像分别是
\[
\ker\varphi=\{g\in G:\varphi(g)=e_H\},
\qquad
\operatorname{im}\varphi=\{\varphi(g):g\in G\}.
\]
核总是正规子群。若 \(N\triangleleft G\) 是正规子群，则左陪集与右陪集一致，商集合 \(G/N\) 上的乘法
\[
(gN)(hN)=ghN
\]
与代表元的选择无关，从而 \(G/N\) 自身成为群。

同态的核精确记录了“被映射压缩掉”的信息。第一同构定理给出
\begin{equation}
G/\ker\varphi\cong \operatorname{im}\varphi.
\label{eq:first-isomorphism-theorem}
\end{equation}
这条结论以后会以许多不同形式出现：商空间、规范冗余、表示的有效作用以及纤维化都可以看作同一种“先识别不可区分对象，再在商对象上工作”的过程。

\begin{derivation}[第一同构定理]
分四步证明：良定义性、同态性、满射性、单射性。整体策略是把“被 \(\varphi\) 等同的元素”打包成陪集，再在陪集上重写映射；四步恰好对应“商对象上函数必须回答的四个问题”。

\emph{第一步：良定义性。}令 \(K=\ker\varphi\)。因为 \(\varphi\) 是同态，\(K\trianglelefteq G\) 是正规子群，故商群 \(G/K\) 良定义，其元素为陪集 \(gK=\{gk:k\in K\}\)。定义
\[
  \widetilde\varphi:G/K\longrightarrow\operatorname{im}\varphi,
  \qquad
  \widetilde\varphi(gK)=\varphi(g).
\]
需验证此定义不依赖陪集代表元选取。若 \(gK=hK\)，则由陪集相等的判据，\(h^{-1}g\in K\)。因 \(K=\ker\varphi\)，
\[
  \varphi(h^{-1}g)=e_H.
\]
由 \(\varphi\) 的同态性，先对逆元有 \(\varphi(h^{-1})=\varphi(h)^{-1}\)（因为 \(\varphi(h^{-1})\varphi(h)=\varphi(h^{-1}h)=\varphi(e)=e_H\)），于是
\[
  \varphi(h^{-1}g)=\varphi(h^{-1})\varphi(g)=\varphi(h)^{-1}\varphi(g)=e_H,
\]
两端左乘 \(\varphi(h)\) 得 \(\varphi(g)=\varphi(h)\)。因此 \(\widetilde\varphi(gK)=\widetilde\varphi(hK)\)，定义良定。注意此处只用到了 \(K\) 的正规性：它保证 \(G/K\) 上的乘法是群运算，而良定义性本身只需要 \(K=\ker\varphi\) 即可。

\emph{第二步：同态性。}对任意 \(gK,hK\in G/K\)，利用 \(G/K\) 中乘法定义 \((gK)(hK)=ghK\)（此定义依赖 \(K\) 正规：若改取代表元 \(gk_1,hk_2\)，则 \((gk_1)(hk_2)=g(k_1h)k_2=g(hk_3)k_2=gh(k_3k_2)\in ghK\)，其中 \(k_3=h^{-1}k_1h\in K\) 用到 \(K\trianglelefteq G\)）。于是
\[
  \widetilde\varphi((gK)(hK))
  =\widetilde\varphi(ghK)
  =\varphi(gh)
  =\varphi(g)\varphi(h)
  =\widetilde\varphi(gK)\,\widetilde\varphi(hK),
\]
第三步用了 \(\varphi\) 的同态性。故 \(\widetilde\varphi\) 是群同态。

\emph{第三步：满射性。}对任意 \(y\in\operatorname{im}\varphi\)，由像的定义存在 \(g\in G\) 使 \(\varphi(g)=y\)。于是 \(\widetilde\varphi(gK)=\varphi(g)=y\)，故 \(\widetilde\varphi\) 满射到 \(\operatorname{im}\varphi\)。注意像取为 \(\operatorname{im}\varphi\) 而非整个 \(H\)，正是满射性的来源；若取 \(H\) 一般不满射。

\emph{第四步：单射性。}设 \(\widetilde\varphi(gK)=e_H\)。由定义 \(\varphi(g)=e_H\)，故 \(g\in\ker\varphi=K\)，因而 \(gK=K\)，即 \(G/K\) 的单位元。核只有单位元，故 \(\widetilde\varphi\) 单射。等价地，\(\ker\widetilde\varphi=\{K\}\) 即 \(G/K\) 的单位元。

综合四步，\(\widetilde\varphi:G/K\xrightarrow{\,\sim\,}\operatorname{im}\varphi\) 是双射同态，即群同构，得到\eqnrefs{eq:first-isomorphism-theorem}。

\emph{举例：循环群与剩余类。}取 \(G=\mathbb Z\)（加法群）、\(H=\mathbb Z_n\)、\(\varphi(k)=k\bmod n\)。则 \(\ker\varphi=n\mathbb Z\)，\(\operatorname{im}\varphi=\mathbb Z_n\)，定理给出 \(\mathbb Z/n\mathbb Z\cong\mathbb Z_n\)。这是模运算严格化的代数来源。

\emph{举例：行列式。}取 \(G=GL_n(\CC)\)、\(H=\CC^\times\)、\(\varphi=\det\)。则 \(\ker\varphi=SL_n(\CC)\)，\(\operatorname{im}\varphi=\CC^\times\)，定理给出
\[
  GL_n(\CC)/SL_n(\CC)\cong\CC^\times.
\]
即“行列式为 \(\lambda\) 的所有矩阵构成一个 \(SL_n\) 陪集”，且这些陪集按 \(\lambda\) 一一对应。

\emph{举例：\(S_3\) 的符号同态。}取 \(G=S_3\)、\(\varphi=\operatorname{sgn}:S_3\to\{\pm1\}\)。则 \(\ker\varphi=A_3\cong\mathbb Z_3\)，\(\operatorname{im}\varphi=\{\pm1\}\)，定理给出 \(S_3/A_3\cong\mathbb Z_2\)。这与 Lagrange 定理一致：\(|S_3|=6=|A_3|\cdot 2\)。

\emph{推广方向。}第一同构定理在群、环、模、Lie 代数、向量空间、范畴论中都有平行版本：环情形给出 \(R/I\cong\operatorname{im}\varphi\)（\(I\) 为理想），模情形给出 \(M/\ker\varphi\cong\operatorname{im}\varphi\)，Lie 代数情形给出 \(\mathfrak g/\ker\varphi\cong\operatorname{im}\varphi\)（\(\ker\varphi\) 为理想）。范畴论中它表现为“每个态射可分解为正则满态射后接单态射”。物理中规范固定正是“在规范等价类构成的商空间上工作”，其代数母结构即此定理。

\emph{物理应用：规范冗余与商空间。}在规范理论中，场空间 \(\mathcal F\) 上的规范群 \(\mathcal G\) 自由作用，物理位形空间是商 \(\mathcal F/\mathcal G\)。第一同构定理保证此商的代数结构良好：若 \(\varphi:\mathcal F\to\mathcal O\) 为映射到可观测量 \(\mathcal O\) 的规范不变量，则 \(\mathcal F/\ker\varphi\cong\operatorname{im}\varphi\)，即可观测量空间与场空间模规范冗余同构。这是 BRST 量子化与规范固定程序的代数基础。在电磁理论中，\(\varphi:A_\mu\mapsto F_{\mu\nu}\) 把规范势映到场强，核为规范变换 \(A_\mu\mapsto A_\mu+\partial_\mu\chi\)，商即物理场空间。此模式在 Yang--Mills、BRST 与引力理论中反复出现，是“先识别不可区分对象，再在商对象上工作”这一普遍策略的代数母结构。在范畴论中，此策略表现为“每个态射分解为正则满态射后接单态射”，是商对象与像对象构造的统一语言。

\emph{总结。}第一同构定理把“同态的核--像”关系精确化为同构，是群论、环论、模论、Lie 代数论与范畴论中商对象构造的代数基础，物理中规范固定与对称约化都依赖此结构。


接下来考察陪集、Lagrange 定理与正规子群。

设 \(H\le G\) 为子群。对任意 \(g\in G\)，左陪集为
\[
gH=\{gh:h\in H\}.
\]
左乘 \(h\mapsto gh\) 是双射，因此每个左陪集与 \(H\) 等势。两个左陪集要么完全相同，要么不相交，于是有限群被等大小的陪集分割：
\begin{equation}
|G|=[G:H]\,|H|.
\label{eq:lagrange-group}
\end{equation}
这里 \([G:H]\) 是 \(H\) 在 \(G\) 中的指数。取 \(H=\langle g\rangle\) 立即得到 \(\operatorname{ord}(g)\mid |G|\)。

正规性可以用共轭作用刻画。子群 \(N\le G\) 正规当且仅当
\[
gNg^{-1}=N,
\qquad \forall g\in G.
\]
这说明正规子群不是任意子对象，而是对群自身的内禀对称性稳定的子对象。

若 \(N\triangleleft G\)、\(H\le G\)，且 \(N\cap H=\{e\}\)、\(NH=G\)，则在共轭作用
\[
\alpha_h(n)=hnh^{-1}
\]
下，\(G\) 可写成半直积 \(N\rtimes_\alpha H\)。半直积与直积的区别正是两个子群之间存在非平凡作用；二面体群
\[
D_n\cong C_n\rtimes C_2
\]
就是最基本的例子，其中 \(C_2\) 通过取逆作用于 \(C_n\)。

接下来考察群作用、轨道与稳定子。

群作用是同态
\[
\Phi:G\longrightarrow \operatorname{Bij}(X),
\]
也可写成 \(g\cdot x\)，满足
\[
e\cdot x=x,
\qquad
(gh)\cdot x=g\cdot(h\cdot x).
\]
给定 \(x\in X\)，其轨道和稳定子分别为
\[
\operatorname{Orb}(x)=\{g\cdot x:g\in G\},
\qquad
G_x=\{g\in G:g\cdot x=x\}.
\]
映射
\[
G/G_x\longrightarrow \operatorname{Orb}(x),
\qquad
gG_x\longmapsto g\cdot x
\]
是双射。因此有限群作用满足 orbit--stabilizer 公式
\begin{equation}
|\operatorname{Orb}(x)|=[G:G_x]
=\frac{|G|}{|G_x|}.
\label{eq:orbit-stabilizer}
\end{equation}

当 \(G\) 作用在自身上且作用取为共轭
\[
g\cdot x=gxg^{-1}
\]
时，轨道就是共轭类，稳定子就是中心化子
\[
C_G(x)=\{g\in G:gx=xg\}.
\]
故共轭类大小为
\[
|\operatorname{Cl}(x)|=[G:C_G(x)].
\]
中心 \(Z(G)\) 中的元素共轭类大小为一。选取所有非中心共轭类代表元 \(x_i\)，得到 class equation
\begin{equation}
|G|=|Z(G)|+\sum_i [G:C_G(x_i)].
\label{eq:class-equation}
\end{equation}
这一公式把群的整体阶数分解为共轭轨道的大小，是有限群结构与 character 理论之间的第一座桥梁。


固定 \(x\in X\)，定义映射
\[
\varphi:G/G_x\longrightarrow\operatorname{Orb}(x),
\qquad
gG_x\longmapsto g\cdot x.
\]
先验证良定义性。若 \(g_1G_x=g_2G_x\)，则 \(g_2^{-1}g_1\in G_x\)，由稳定子定义 \((g_2^{-1}g_1)\cdot x=x\)，因此
\[
g_1\cdot x
=g_2\cdot\bigl((g_2^{-1}g_1)\cdot x\bigr)
=g_2\cdot x,
\]
即 \(\varphi(g_1G_x)=\varphi(g_2G_x)\)。再证单射。若 \(g_1\cdot x=g_2\cdot x\)，两端左乘 \(g_2^{-1}\) 得
\[
(g_2^{-1}g_1)\cdot x=x,
\]
故 \(g_2^{-1}g_1\in G_x\)，即 \(g_1G_x=g_2G_x\)。满射是直接的：\(\operatorname{Orb}(x)\) 中每个元素形如 \(g\cdot x\)，恰为 \(\varphi(gG_x)\) 的像。因此 \(\varphi\) 是双射，
\[
G/G_x\simeq \operatorname{Orb}(x).
\]
对有限群取基数，\(|G/G_x|=[G:G_x]=|G|/|G_x|\)，从而得到\eqnrefs{eq:orbit-stabilizer}。

对共轭作用取 \(X=G\)，即 \(g\cdot x=gxg^{-1}\)。此时稳定子
\[
G_x=\{g\in G:gxg^{-1}=x\}=\{g\in G:gx=xg\}=C_G(x)
\]
恰为中心化子，轨道 \(\operatorname{Orb}(x)=\{gxg^{-1}:g\in G\}=\operatorname{Cl}(x)\) 恰为共轭类。每个元素恰属于一个共轭类，故群可写成共轭类的不交并
\[
G=Z(G)\;\sqcup\;\bigsqcup_i\operatorname{Cl}(x_i),
\]
其中 \(x_i\) 遍取非中心共轭类代表元。中心元素 \(z\in Z(G)\) 的稳定子为整个 \(G\)，由 orbit--stabilizer 公式其轨道大小 \(|G|/|G|=1\)；非中心元素 \(x_i\) 的轨道大小为 \([G:C_G(x_i)]\)。对不交并取基数，
\[
|G|=|Z(G)|+\sum_i [G:C_G(x_i)],
\]
即\eqnrefs{eq:class-equation}。

\emph{举例：\(S_3\) 的 class equation。}\(S_3\) 有三个共轭类：\(\{e\}\)、\(\{(12),(13),(23)\}\)、\(\{(123),(132)\}\)。中心 \(Z(S_3)=\{e\}\)，故
\[
6=1+\underbrace{3}_{[(12)\text{ 类}]}+\underbrace{2}_{[(123)\text{ 类}]}.
\]
这里 \([S_3:C_{S_3}((12))]=3\)（\(C_{S_3}((12))=\{e,(12)\}\cong\mathbb Z_2\)，\([S_3:\mathbb Z_2]=3\)），\([S_3:C_{S_3}((123))]=2\)（\(C_{S_3}((123))=\{e,(123),(132)\}\cong\mathbb Z_3\)，\([S_3:\mathbb Z_3]=2\)）。

\emph{举例：\(D_4\) 的 class equation。}二面体群 \(D_4=\langle r,s:r^4=s^2=e,\,srs=r^{-1}\rangle\) 有五个共轭类：\(\{e\}\)、\(\{r^2\}\)、\(\{r,r^3\}\)、\(\{s,sr^2\}\)、\(\{sr,sr^3\}\)。中心 \(Z(D_4)=\{e,r^2\}\)，故
\[
8=2+\underbrace{2}_{[r\text{ 类}]}+\underbrace{2}_{[s\text{ 类}]}+\underbrace{2}_{[sr\text{ 类}]}.
\]
注意每个非中心类大小均为 \(2\)，恰为 \([D_4:C_{D_4}(x)]\)。

\emph{应用：\(p\)-群有非平凡中心。}若 \(|G|=p^k\)（\(p\) 素），class equation 中每个 \([G:C_G(x_i)]\) 整除 \(|G|\) 且大于 \(1\)，故为 \(p\) 的幂。于是
\[
|Z(G)|=|G|-\sum_i [G:C_G(x_i)]
\]
右端每项被 \(p\) 整除，故 \(p\mid |Z(G)|\)。这推出 \(p\)-群必有非平凡中心，是 \(p\)-群结构定理与 Sylow 定理的起点。

\emph{推广方向。}对紧致 Lie 群 \(G\)，class equation 升级为 Weyl 积分公式：对类函数 \(f\)
\[
\int_G f(g)\,\dd g=\frac{1}{|W|}\int_T f(t)\,|\Delta(t)|^2\,\dd t,
\]
其中 \(T\) 为 maximal torus，\(W\) 为 Weyl 群，\(\Delta\) 为 Weyl 分母。有限群的 class equation 可视为 Lie 群情形“离散化”的特例。物理中，class equation 是把表示论约化到共轭类（character 在共轭类上常值）的组合基础，Burnside 引理与 Frobenius 公式都依赖它。



接下来考察Burnside 轨道计数、$p$-群与群扩张。

Orbit--stabilizer 描述单个轨道；若要直接计算轨道总数，可以把所有固定点信息求平均。设有限群 $G$ 作用在有限集合 $X$ 上，记
\[
X^g:=\{x\in X:g\cdot x=x\}.
\]
Burnside 引理给出
\begin{equation}
\boxed{
|X/G|
=\frac1{|G|}\sum_{g\in G}|X^g|
}.
\label{eq:burnside-orbit-counting}
\end{equation}
它把“按轨道分类”的问题转化为“逐个群元素数固定点”。Polya 计数、晶格着色和组合对称性都由这一式继续发展。


\emph{思路：}对"群元素固定样本点"的二元关系做双重计数——先按群元素计数，再按样本点计数，后者按轨道分组后每轨道贡献恰为 $|G|$。

\emph{第一步：构造计数集合。}定义
\begin{equation*}
 \mathcal S
 :=\{(g,x)\in G\times X:g\cdot x=x\},
\end{equation*}
即所有"群元素 $g$ 固定点 $x$"的配对。

\emph{第二步：按群元素计数。}固定 $g\in G$，满足 $g\cdot x=x$ 的 $x$ 恰是 $g$ 的不动点集 $X^g$。因此
\begin{equation*}
 |\mathcal S|
 =\sum_{g\in G}|X^g|.
\end{equation*}

\emph{第三步：按样本点计数。}固定 $x\in X$，满足 $g\cdot x=x$ 的 $g$ 恰是 $x$ 的稳定子 $G_x$。因此
\begin{equation*}
 |\mathcal S|
 =\sum_{x\in X}|G_x|.
\end{equation*}

\emph{第四步：按轨道分组。}将上式按轨道 $\mathcal O\in X/G$ 分组。同一轨道内各点的稳定子彼此共轭，故大小相同。由 orbit--stabilizer 定理，
\begin{equation*}
 |\mathcal O|\,|G_x|=|G|,
 \qquad x\in\mathcal O.
\end{equation*}
因此每个轨道对第三步总和的贡献为
\begin{equation*}
 \sum_{x\in\mathcal O}|G_x|
 =|\mathcal O|\,|G_x|
 =|G|.
\end{equation*}
轨道数为 $|X/G|$，故
\begin{equation*}
 |\mathcal S|=|G|\,|X/G|.
\end{equation*}

\emph{第五步：合并。}将第二步与第四步结果等同：
\begin{equation*}
 \sum_{g\in G}|X^g|=|G|\,|X/G|,
\end{equation*}
两边除以 $|G|$ 即得\eqnrefs{eq:burnside-orbit-counting}。


Class equation 对 $p$-群尤其强。若 $|G|=p^n$，每个非中心共轭类大小
\[
[G:C_G(x)]
\]
既是 $p$ 的幂又大于 $1$，因此都被 $p$ 整除。由 class equation，
\[
|G|
=|Z(G)|+
\sum_i[G:C_G(x_i)]
\]
两边模 $p$，得到
\[
|Z(G)|\equiv0\pmod p.
\]
所以有限 $p$-群的中心一定非平凡。特别地，若 $|G|=p^2$，则 $|Z(G)|=p$ 或 $p^2$。若为 $p^2$，群 Abel；若为 $p$，则 $G/Z(G)$ 阶为 $p$，故循环，而“商去中心后循环”会迫使原群 Abel，矛盾。因此
\[
\boxed{|G|=p^2\quad\Longrightarrow\quad G\text{ Abel}.}
\]
这类结论显示 class equation 不只是 character 前的准备，而是有限群结构本身的强约束。



$p$-群中心定理进一步通向有限群中最重要的存在--共轭定理之一。先证明 Cauchy 定理：若素数 $p$ 整除 $|G|$，则 $G$ 中存在阶为 $p$ 的元素。考虑
\[
\mathcal X:=\{(g_1,\ldots,g_p)\in G^p:g_1g_2\cdots g_p=e\}.
\]
前 $p-1$ 个分量任取后最后一个被唯一确定，因此
\[
|\mathcal X|=|G|^{p-1}\equiv0\pmod p.
\]
循环群 $C_p$ 通过循环移位作用在 $\mathcal X$ 上。每个非固定轨道大小都是 $p$，故固定点数与 $|\mathcal X|$ 模 $p$ 同余。固定点恰为
\[
(g,g,\ldots,g),\qquad g^p=e.
\]
恒等元给出一个固定点，而固定点总数被 $p$ 整除，所以还存在 $g\neq e$ 满足 $g^p=e$；由 $p$ 为素数，$g$ 的阶正是 $p$。

\textbf{定理（Sylow 定理）.}\enspace \label{thm:math-group-sylow}
设
\[
|G|=p^n m,
\qquad p\nmid m.
\]
则：
\begin{enumerate}[label=(\roman*)]
\item 存在阶为 $p^n$ 的子群 $P\le G$，称为 Sylow $p$-subgroup；
\item 任意两个 Sylow $p$-subgroups 共轭；
\item 若 $n_p$ 是 Sylow $p$-subgroups 的个数，则
\begin{equation}
\boxed{
n_p\equiv1\pmod p,
\qquad
n_p\mid m.
}
\label{eq:math-group-sylow-count}
\end{equation}
\end{enumerate}


先证存在性，对 $|G|$ 归纳。若 $p\mid|Z(G)|$，Cauchy 定理给出中心子群
\[
C\le Z(G),\qquad |C|=p.
\]
商群 $G/C$ 的阶为 $p^{n-1}m$。由归纳假设，它有阶 $p^{n-1}$ 的子群 $\overline P$；其在商映射下的原像 $P$ 满足
\[
|P|=|C|\,|\overline P|=p^n.
\]

若 $p\nmid|Z(G)|$，由 class equation
\[
|G|=|Z(G)|+\sum_i[G:C_G(x_i)]
\]
可知至少有一个非中心元素 $x_i$ 使
\[
p\nmid[G:C_G(x_i)].
\]
因此 $|C_G(x_i)|$ 含有完整因子 $p^n$。又因 $x_i$ 非中心，$C_G(x_i)$ 是真子群；对它使用归纳假设，即得到阶 $p^n$ 的子群。

再证共轭性。设 $P,Q$ 都是 Sylow $p$-subgroups。让 $P$ 左作用于陪集空间 $G/Q$。因为
\[
|G/Q|=m
\]
不被 $p$ 整除，而任意非平凡 $P$-轨道大小都是 $p$ 的正幂，所以必有固定陪集 $gQ$。固定意味着
\[
PgQ=gQ
\quad\Longleftrightarrow\quad
P\subseteq gQg^{-1}.
\]
两边阶都为 $p^n$，故
\[
P=gQg^{-1}.
\]

最后，所有 Sylow $p$-subgroups 构成一个共轭轨道，所以
\[
n_p=[G:N_G(P)].
\]
因为 $P\subseteq N_G(P)$，指标 $n_p$ 必整除 $[G:P]=m$。另一方面让 $P$ 通过共轭作用在 Sylow 集合上。$P$ 自身是固定点；若另一个 Sylow $Q$ 也被整个 $P$ 固定，则 $P\subseteq N_G(Q)$，从而 $PQ$ 是 $p$-subgroup。Sylow 极大性迫使 $P=Q$。因此除 $P$ 外所有轨道大小都被 $p$ 整除，于是
\[
n_p\equiv1\pmod p.
\]
这得到\eqnrefs{eq:math-group-sylow-count}。
\end{derivation}

Sylow 计数常常无需构造群就能排除大量可能。例如若
\[
|G|=pq,
\qquad p<q,
\]
其中 $p,q$ 为素数，则 Sylow $q$-subgroup 个数满足
\[
n_q\mid p,
\qquad
n_q\equiv1\pmod q.
\]
由于 $p<q$，只能有 $n_q=1$，故阶 $q$ 子群自动正规。有限群的结构分析由此形成一条清楚链条：Lagrange 限制可能阶数，Cauchy 保证素数阶元素存在，Sylow 再控制最大 $p$-子群的存在、共轭与正规性。

群扩张把“由正规子群和商群重建整个群”系统化。一个短正合列
\begin{equation}
1\longrightarrow N
\xrightarrow{\iota}G
\xrightarrow{\pi}Q
\longrightarrow1
\label{eq:group-extension-short-exact}
\end{equation}
意味着 $\iota(N)=\ker\pi$，且 $Q\cong G/N$。若存在群同态截面
\[
s:Q\to G,
\qquad
\pi\circ s=\IdMap_Q,
\]
则扩张称为 split，并有
\[
G\cong N\rtimes Q.
\]
所以半直积并不是一种偶然乘法，而是“短正合列存在同态截面”的精确结构。

若只任选集合截面 $s$ 而不要求它保持乘法，则失配由 factor set
\[
c(q_1,q_2)
:=s(q_1)s(q_2)s(q_1q_2)^{-1}
\in N
\]
测量。若 $N$ Abel，且 $Q$ 通过共轭诱导作用 $q\cdot n$ 于 $N$，结合律迫使
\begin{equation}
c(q_1,q_2)
+c(q_1q_2,q_3)
=q_1\cdot c(q_2,q_3)
+c(q_1,q_2q_3),
\label{eq:group-extension-2cocycle}
\end{equation}
这里把 $N$ 的群运算写成加法。这正是 $2$-cocycle 条件。改变截面
\[
s'(q)=b(q)s(q)
\]
会把 $c$ 改变一个 coboundary；因此在固定 $Q$-作用下，Abelian 扩张的等价类由群上同调 $H^2(Q,N)$ 组织。这里不展开完整群上同调，但这已经给出一个重要层级：
\[
\text{直积}
\subset
\text{半直积}
\subset
\text{一般扩张},
\]
而“不 split 的程度”由 cocycle 数据记录。

\section{有限群的线性表示}
\label{sec:finite-group-representations}

接下来考察表示、intertwiner 与不可约性。

设 \(V\) 是有限维复向量空间。群 \(G\) 在 \(V\) 上的线性表示是同态
\[
\rho:G\longrightarrow GL(V).
\]
若子空间 \(W\subset V\) 对所有 \(\rho(g)\) 都不变，即
\[
\rho(g)W\subset W,
\qquad \forall g\in G,
\]
则称 \(W\) 为不变子空间。若除 \(\{0\}\) 与 \(V\) 外不存在非平凡不变子空间，则 \(\rho\) 不可约。

两个表示 \((V,\rho)\)、\((W,\sigma)\) 之间满足
\[
A\rho(g)=\sigma(g)A,
\qquad \forall g\in G
\]
的线性映射 \(A:V\to W\) 称为 intertwiner。Schur 引理说明，不可约表示之间的 intertwiner 极其受限：若 \(V,W\) 不同构，则 \(A=0\)；若 \(V=W\)、\(\rho=\sigma\)，则在复数域上
\begin{equation}
A=\lambda I_V.
\label{eq:schur-scalar}
\end{equation}

有限群的任意复表示都可以酉化。给定任意正定 Hermitian 内积 \((\cdot,\cdot)_0\)，定义群平均
\[
\langle v,w\rangle_G
=\frac1{|G|}\sum_{g\in G}
(\rho(g)v,\rho(g)w)_0.
\]
则 \(\langle\cdot,\cdot\rangle_G\) 正定且满足
\[
\langle\rho(h)v,\rho(h)w\rangle_G
=\langle v,w\rangle_G.
\]
因此以后讨论有限群 character 时可不失一般性地假设表示是酉表示。

\begin{derivation}[Schur 引理与 Maschke 完全可约性]
\emph{策略：}Schur 引理先利用「核与像是不变子空间」得出非零 intertwiner 必为同构，再借本征值把它压成标量；Maschke 定理则把一个任意投影做群平均，得到既与表示对易、又以 \(W\) 为像的等变投影，从而切出不变补子空间。

\emph{第一步：核与像都是不变子空间。}设 \(A:V\to W\) 是不可约表示之间的 intertwiner，即
\[
A\rho(g)=\sigma(g)A.
\]
若 \(v\in\ker A\)，则
\begin{equation*}
  A(\rho(g)v)=\sigma(g)A(v)=\sigma(g)\,0=0,
\end{equation*}
故 \(\rho(g)v\in\ker A\)，即 \(\ker A\subset V\) 是 \(\rho\)-不变子空间。若 \(w=A(v)\in\operatorname{im}A\)，则
\begin{equation*}
  \sigma(g)w=\sigma(g)A(v)=A(\rho(g)v)\in\operatorname{im}A,
\end{equation*}
故 \(\operatorname{im}A\subset W\) 是 \(\sigma\)-不变子空间。

\emph{第二步：不可约性迫使它们平凡。}由 \(\rho,\sigma\) 不可约，\(V\) 的不变子空间只有 \(\{0\},V\)，\(W\) 的只有 \(\{0\},W\)，故
\[
\ker A=\{0\}\ \text{或}\ V,
\qquad
\operatorname{im}A=\{0\}\ \text{或}\ W.
\]
若 \(A\neq0\)，则 \(\ker A\neq V\)（否则 \(A\) 处处为零），故 \(\ker A=\{0\}\) 即单射；同时 \(\operatorname{im}A\neq\{0\}\)，故 \(\operatorname{im}A=W\) 即满射。于是非零 \(A\) 必为同构。

\emph{第三步：标量情形。}若进一步 \(V=W\)、\(\rho=\sigma\)，且 \(V\) 是有限维复空间，则 \(A\) 的特征多项式在 \(\mathbb C\) 上有根，取本征值 \(\lambda\)，算子
\[
A-\lambda I
\]
仍是 intertwiner：
\begin{equation*}
  (A-\lambda I)\rho(g)
  =A\rho(g)-\lambda\rho(g)
  =\rho(g)A-\rho(g)\lambda
  =\rho(g)(A-\lambda I).
\end{equation*}
它又有非零核 \(\ker(A-\lambda I)\neq\{0\}\)（\(\lambda\) 是本征值），故不是同构。由第二步非零 intertwiner 必为同构，只能 \(A-\lambda I=0\)，即 \(A=\lambda I\)，这就是\eqnrefs{eq:schur-scalar}。

\emph{第四步：Maschke 的等变投影。}设 \(W\subset V\) 是不变子空间，先任取一个线性投影 \(P_0:V\to W\)（即 \(P_0^2=P_0\)、\(\operatorname{im}P_0=W\)）。作群平均
\[
P=\frac1{|G|}\sum_{g\in G}
\rho(g)P_0\rho(g)^{-1}.
\]
对任意 \(h\in G\)，换元 \(g'=hg\) 得
\[
\rho(h)P\rho(h)^{-1}
=\frac1{|G|}\sum_{g'\in G}
\rho(g')P_0\rho(g')^{-1}
=P,
\]
所以 \(P\) 与表示对易。

\emph{第五步：验证 \(P\) 是要的投影。}因 \(W\) 不变，对 \(w\in W\) 有
\[
\rho(g)^{-1}w\in W,
\qquad
P_0\rho(g)^{-1}w=\rho(g)^{-1}w,
\]
故 \(\rho(g)P_0\rho(g)^{-1}w=w\)，平均得 \(Pw=w\)。由此：\(W\subset\operatorname{im}P\)；反向每个 \(\rho(g)P_0\rho(g)^{-1}v\) 落在 \(W\)（因 \(P_0\) 的像在 \(W\) 且 \(\rho(g)W\subset W\)），故 \(P v\in W\)，即 \(\operatorname{im}P\subset W\)，合得 \(\operatorname{im}P=W\)。再对任意 \(v\)，\(P v\in W\)，用 \(Pw=w\) 得 \(P(P v)=P v\)，即 \(P^2=P\)。最后 \(\ker P\) 也是不变子空间：若 \(v\in\ker P\)，则
\begin{equation*}
  P(\rho(h)v)=\rho(h)P v=0,
\end{equation*}
故 \(\rho(h)v\in\ker P\)。于是 \(P\) 是投影，且像与核都是不变子空间，线性代数给出
\[
V=\operatorname{im}P\oplus\ker P=W\oplus\ker P.
\]

\emph{第六步：归纳完成。}若 \(W\) 或 \(\ker P\) 仍可约，就继续在其内部重复上述切分；有限维保证过程在有限步后停止于不可约子空间的直和，即有限群复表示的完全可约性。这一结论只在特征与 \(|G|\) 互素时成立，特别是 \(\mathbb C\) 上恒成立，是有限群表示论能与对角化类比的根本原因。


接下来考察character、正交关系与不可约重数。

表示 \(\rho\) 的 character 定义为
\[
\chi_\rho(g)=\operatorname{tr}\rho(g).
\]
因为迹在相似变换下不变，\(\chi_\rho\) 在每个共轭类上为常数。若
\[
V\cong \bigoplus_\alpha m_\alpha V_\alpha
\]
是不可约分解，则
\[
\chi_V=\sum_\alpha m_\alpha\chi_\alpha.
\]

在所有复值 class function 上定义内积
\begin{equation}
\langle f,h\rangle_G
=\frac1{|G|}\sum_{g\in G}
\overline{f(g)}h(g).
\label{eq:character-inner-product}
\end{equation}
有限群不可约 character 满足
\begin{equation}
\langle\chi_\alpha,\chi_\beta\rangle_G
=\delta_{\alpha\beta}.
\label{eq:character-orthogonality}
\end{equation}
因此不可约表示 \(V_\alpha\) 在 \(V\) 中的重数为
\begin{equation}
m_\alpha
=\langle\chi_\alpha,\chi_V\rangle_G.
\label{eq:character-multiplicity}
\end{equation}
这比直接寻找不变子空间有效得多：矩阵表示的分解问题被化成有限个迹的求和。


设 \(\rho_\alpha:G\to U(V_\alpha)\)、\(\rho_\beta:G\to U(V_\beta)\) 是不可约酉表示。对任意线性映射 \(A:V_\alpha\to V_\beta\)，定义群平均
\[
\mathcal A(A)
=\frac1{|G|}\sum_{g\in G}
\rho_\beta(g)A\rho_\alpha(g)^{-1}.
\]
直接换元可得
\[
\rho_\beta(h)\mathcal A(A)
=\mathcal A(A)\rho_\alpha(h),
\]
因此 \(\mathcal A(A)\) 是 intertwiner。若 \(\alpha\neq\beta\)，Schur 引理给出 \(\mathcal A(A)=0\)。若 \(\alpha=\beta\)，则
\[
\mathcal A(A)=c(A)I.
\]
取迹并利用迹的循环性，得到
\[
\operatorname{tr}\mathcal A(A)
=\operatorname{tr}A
=c(A)d_\alpha,
\]
故
\[
\mathcal A(A)=\frac{\operatorname{tr}A}{d_\alpha}I.
\]

取矩阵单位 \(A=E_{ji}\)，比较矩阵元，可得矩阵元正交关系
\[
\frac1{|G|}\sum_{g\in G}
[\rho_\alpha(g)]_{ij}
\overline{[\rho_\beta(g)]_{k\ell}}
=
\frac{\delta_{\alpha\beta}}{d_\alpha}
\delta_{ik}\delta_{j\ell}.
\]
令 \(i=j\)、\(k=\ell\) 并对 \(i,k\) 求和，便得到\eqnrefs{eq:character-orthogonality}。再把
\[
\chi_V=\sum_\alpha m_\alpha\chi_\alpha
\]
与 \(\chi_\beta\) 作内积，即得到\eqnrefs{eq:character-multiplicity}。


接下来考察投影算符与正则表示。

character 正交关系可以直接构造等型分量的投影。设 \(\rho:G\to GL(V)\)，\(\alpha\) 是维数为 \(d_\alpha\) 的不可约表示，则
\begin{equation}
P_\alpha
=
\frac{d_\alpha}{|G|}
\sum_{g\in G}
\chi_\alpha(g^{-1})\rho(g)
\label{eq:character-projector}
\end{equation}
投影到 \(V\) 中所有与 \(V_\alpha\) 同构的不可约副本之和。它满足
\[
P_\alpha^2=P_\alpha,
\qquad
P_\alpha P_\beta=0\quad(\alpha\neq\beta),
\qquad
\sum_\alpha P_\alpha=I.
\]
在分子振动、晶体正常模和简并本征空间的对称性分解中，真正使用的正是\eqnrefs{eq:character-projector}，而不是单纯的 character 表。

群在函数空间 \(\mathbb C[G]\) 上的左正则表示定义为
\[
(L_hf)(g)=f(h^{-1}g).
\]
其维数为 \(|G|\)。正则表示包含每个不可约表示 \(V_\alpha\) 恰好 \(d_\alpha\) 次，因此
\begin{equation}
|G|=\sum_\alpha d_\alpha^2.
\label{eq:sum-irrep-dimensions}
\end{equation}
这条恒等式既是不可约表示分类的强约束，也是 finite Fourier transform 的维数守恒式。


由 \(\chi_\alpha\) 是 class function，\eqnrefs{eq:character-projector} 与每个 \(\rho(h)\) 对易。直接计算
\begin{align*}
\rho(h)P_\alpha\rho(h)^{-1}
&=\frac{d_\alpha}{|G|}
\sum_{g\in G}
\chi_\alpha(g^{-1})\rho(h)\rho(g)\rho(h)^{-1}\\
&=\frac{d_\alpha}{|G|}
\sum_{g\in G}
\chi_\alpha(g^{-1})\rho(hgh^{-1}).
\end{align*}
令 \(g'=hgh^{-1}\)。由共轭是 \(G\) 上双射，求和指标可换为 \(g'\)；又由 \(\chi_\alpha\) 是 class function，\(\chi_\alpha(g^{-1})=\chi_\alpha(g'^{-1})\)，因此
\[
\rho(h)P_\alpha\rho(h)^{-1}
=\frac{d_\alpha}{|G|}
\sum_{g'\in G}
\chi_\alpha(g'^{-1})\rho(g')
=P_\alpha.
\]
故 \(P_\alpha\) 是 \(G\)-等变算子。将 \(V\) 分解为不可约子空间直和 \(V=\bigoplus_\beta V_\beta\)，在每个 \(V_\beta\) 上 \(P_\alpha\) 与不可约表示 \(\rho|_{V_\beta}\) 交换。由 Schur 引理，\(P_\alpha|_{V_\beta}=c_{\alpha\beta}I_{V_\beta}\)。对该等式两边取迹，左端为
\begin{align*}
\operatorname{tr}_{V_\beta}(P_\alpha|_{V_\beta})
&=\operatorname{tr}_{V_\beta}\!\left(
\frac{d_\alpha}{|G|}
\sum_{g\in G}
\chi_\alpha(g^{-1})\rho_\beta(g)\right)\\
&=\frac{d_\alpha}{|G|}
\sum_{g\in G}
\chi_\alpha(g^{-1})\operatorname{tr}\rho_\beta(g)\\
&=\frac{d_\alpha}{|G|}
\sum_{g\in G}
\chi_\alpha(g^{-1})\chi_\beta(g),
\end{align*}
右端为 \(c_{\alpha\beta}d_\beta\)，故
\[
c_{\alpha\beta}d_\beta
=
\frac{d_\alpha}{|G|}
\sum_{g\in G}
\chi_\alpha(g^{-1})\chi_\beta(g).
\]
对酉表示有 \(\rho_\alpha(g^{-1})=\rho_\alpha(g)^{-1}=\rho_\alpha(g)^\dagger\)，因此
\[
\chi_\alpha(g^{-1})
=\operatorname{tr}\rho_\alpha(g^{-1})
=\overline{\operatorname{tr}\rho_\alpha(g)}
=\overline{\chi_\alpha(g)}.
\]
代入并使用\eqnrefs{eq:character-orthogonality}，
\[
c_{\alpha\beta}d_\beta
=\frac{d_\alpha}{|G|}
\sum_{g\in G}
\overline{\chi_\alpha(g)}\chi_\beta(g)
=d_\alpha\delta_{\alpha\beta}.
\]
当 \(\alpha=\beta\) 时 \(d_\alpha=d_\beta\neq0\)，故 \(c_{\alpha\alpha}=1\)；当 \(\alpha\neq\beta\) 时 \(\delta_{\alpha\beta}=0\)，故 \(c_{\alpha\beta}=0\)。于是 \(P_\alpha|_{V_\beta}=\delta_{\alpha\beta}I_{V_\beta}\)，即 \(P_\alpha\) 在 \(\alpha\)-等型分量上为恒等、其余分量上为零。因此 \(P_\alpha\) 恰为 \(\alpha\)-等型分量的投影。
\end{derivation}

\section{张量积、Clebsch--Gordan 分解与对称性选择}
\label{sec:tensor-product-representations}

接下来考察张量积表示与不可约重数。

若 \((V,\rho)\)、\((W,\sigma)\) 是 \(G\) 的表示，则张量积空间 \(V\otimes W\) 上自然有表示
\[
(\rho\otimes\sigma)(g)
=\rho(g)\otimes\sigma(g).
\]
其 character 为
\begin{equation}
\chi_{V\otimes W}(g)
=\chi_V(g)\chi_W(g).
\label{eq:tensor-character-product}
\end{equation}
因此若
\[
V_\alpha\otimes V_\beta
\cong
\bigoplus_\gamma
N_{\alpha\beta}^{\ \ \gamma}V_\gamma,
\]
则 Clebsch--Gordan 重数由
\begin{equation}
N_{\alpha\beta}^{\ \ \gamma}
=
\frac1{|G|}
\sum_{g\in G}
\overline{\chi_\gamma(g)}
\chi_\alpha(g)\chi_\beta(g)
\label{eq:cg-character-multiplicity}
\end{equation}
给出。所谓“两个角动量如何耦合”或“两个模式的直积含有哪些对称类型”，在抽象层面都是同一个张量积表示分解问题。

线性算子空间也可以写成张量积。若 \(V\) 是有限维复表示，则
\[
\operatorname{End}(V)
\cong V\otimes V^*.
\]
群通过
\[
A\longmapsto \rho(g)A\rho(g)^{-1}
\]
作用在 \(\operatorname{End}(V)\) 上。一个算符是否允许连接两个不可约子空间，取决于相应张量积中是否包含平凡表示。这个表述是选择定则的表示论母结构；具体到旋转群时才变成 Wigner--Eckart 定理和熟悉的 \(3j\) 符号条件。

接下来考察不变量与 Hom 空间。

表示之间的 intertwiner 空间可写成
\[
\operatorname{Hom}_G(V,W)
=
\{A\in\operatorname{Hom}(V,W):A\rho(g)=\sigma(g)A\}.
\]
另一方面，利用
\[
\operatorname{Hom}(V,W)\cong V^*\otimes W,
\]
有
\begin{equation}
\operatorname{Hom}_G(V,W)
\cong (V^*\otimes W)^G,
\label{eq:intertwiner-invariants}
\end{equation}
其中右边表示群不变向量空间。于是“求所有对称性允许的线性映射”与“求张量积表示中的 singlet”完全等价。

若 \(V_\alpha,V_\beta,V_\gamma\) 不可约，则三线性不变量的数目为
\[
\dim (V_\alpha\otimes V_\beta\otimes V_\gamma)^G.
\]
这解释了为什么耦合系数不是额外添加的经验常数，而是表示空间中不变张量的坐标。

\begin{derivation}[张量积 character 与耦合重数]
取 \(V\) 与 \(W\) 的基，使 \(\rho(g)\)、\(\sigma(g)\) 分别表示为矩阵。具体地，设 \(V\) 基为 \(\{e_i\}_{i=1}^n\)，\(W\) 基为 \(\{f_a\}_{a=1}^m\)，则 \(V\otimes W\) 的基为 \(\{e_i\otimes f_a\}_{i,a}\)，且
\[
(\rho(g)\otimes\sigma(g))(e_i\otimes f_a)
=\rho(g)e_i\otimes\sigma(g)f_a
=\sum_{j,b}\rho_{ji}(g)\sigma_{ba}(g)\,e_j\otimes f_b.
\]
张量积矩阵的迹满足
\[
\operatorname{tr}(A\otimes B)
=\sum_{i,a}(A\otimes B)_{(i,a),(i,a)}
=\sum_{i,a}A_{ii}B_{aa}
=\Bigl(\sum_i A_{ii}\Bigr)\Bigl(\sum_a B_{aa}\Bigr)
=\operatorname{tr}A\,\operatorname{tr}B,
\]
故
\[
\chi_{V\otimes W}(g)
=
\operatorname{tr}(\rho(g)\otimes\sigma(g))
=
\chi_V(g)\chi_W(g),
\]
得到\eqnrefs{eq:tensor-character-product}。

若
\[
V_\alpha\otimes V_\beta
\cong
\bigoplus_\gamma N_{\alpha\beta}^{\ \ \gamma}V_\gamma,
\]
则取 character 后
\[
\chi_\alpha\chi_\beta
=
\sum_\gamma N_{\alpha\beta}^{\ \ \gamma}\chi_\gamma.
\]
与 \(\overline{\chi_\delta}\) 按\eqnrefs{eq:character-inner-product} 作内积并使用不可约 character 正交性 \(\langle\chi_\gamma,\chi_\delta\rangle_G=\delta_{\gamma\delta}\)，
\[
\langle\chi_\delta,\chi_\alpha\chi_\beta\rangle_G
=\sum_\gamma N_{\alpha\beta}^{\ \ \gamma}\langle\chi_\delta,\chi_\gamma\rangle_G
=\sum_\gamma N_{\alpha\beta}^{\ \ \gamma}\delta_{\delta\gamma}
=N_{\alpha\beta}^{\ \ \delta},
\]
即\eqnrefs{eq:cg-character-multiplicity}。

\emph{举例：\(S_3\) 的二维表示自乘。}\(S_3\) 有三个不可约表示：平凡 \(\chi_1\)、符号 \(\chi_\mathrm{sgn}\)、二维标准 \(\chi_2\)。其 character 表为
\[
\begin{array}{c|ccc}
 & e & (12) & (123)\\
\hline
\chi_1 & 1 & 1 & 1\\
\chi_\mathrm{sgn} & 1 & -1 & 1\\
\chi_2 & 2 & 0 & -1
\end{array}
\]
计算 \(\chi_2\chi_2\)：在 \(e,(12),(123)\) 处取值 \(4,0,1\)。重数为
\[
N_{22}^{\ \ 1}=\langle\chi_1,\chi_2\chi_2\rangle=\tfrac16(1\cdot4+3\cdot0+2\cdot1)=1,
\]
\[
N_{22}^{\ \ \mathrm{sgn}}=\langle\chi_\mathrm{sgn},\chi_2\chi_2\rangle=\tfrac16(1\cdot4+3\cdot0+2\cdot1)=1,
\]
\[
N_{22}^{\ \ 2}=\langle\chi_2,\chi_2\chi_2\rangle=\tfrac16(2\cdot4+0\cdot0+(-1)\cdot2\cdot1)=1.
\]
故 \(V_2\otimes V_2\cong V_1\oplus V_\mathrm{sgn}\oplus V_2\)，即二维表示自乘后含每个不可约表示各一次。

\emph{物理应用：角动量耦合。}对 \(SU(2)\)，不可约表示由自旋 \(j\in\frac12\mathbb Z_{\ge0}\) 标记，\(\dim V_j=2j+1\)。Clebsch--Gordan 级数
\[
V_{j_1}\otimes V_{j_2}\cong\bigoplus_{j=|j_1-j_2|}^{j_1+j_2}V_j
\]
正是\eqnrefs{eq:cg-character-multiplicity} 在 \(SU(2)\) 情形的具体化。例如 \(V_{1/2}\otimes V_{1/2}\cong V_0\oplus V_1\)（自旋单态 + 三重态），\(V_1\otimes V_{1/2}\cong V_{1/2}\oplus V_{3/2}\)。物理中两个角动量耦合后的总角动量本征态分解即此公式的实现，Clebsch--Gordan 系数是相应基变换的矩阵元。

\emph{推广方向。}对紧致 Lie 群 \(G\)，公式升级为
\[
N_{\lambda\mu}^{\ \ \nu}
=\int_G \overline{\chi_\nu(g)}\chi_\lambda(g)\chi_\mu(g)\,\dd g,
\]
其中 \(\dd g\) 为 Haar 测度。对一般紧群，Littlewood--Richardson 规则给出 \(U(n)\) 情形的组合算法；对 \(SU(2)\) 即 Clebsch--Gordan；对 \(SU(3)\) 即夸克味道对称中的重子--介子分解。在算子代数与共形场论中，同一公式以 fusion rule \(N_{ab}^{\ \ c}\) 形式出现，是 Verlinde 公式的代数核心。


接下来考察诱导表示与 Frobenius reciprocity。

张量积从已有表示构造新表示；另一种同样基本的构造是从子群表示向整个群提升。设 $H\le G$，$W$ 是 $H$-表示，作用记为 $\sigma$。有限群情形可定义
\[
\operatorname{Ind}_H^G W
:=\left\{
 f:G\to W:
 f(gh)=\sigma(h^{-1})f(g)
\right\}.
\]
$G$ 通过左平移作用：
\[
(g_0\cdot f)(g)
:=f(g_0^{-1}g).
\]
若选左陪集代表元 $r_1,\ldots,r_m$，其中 $m=[G:H]$，则 $f$ 由 $f(r_a)$ 唯一决定，所以
\[
\dim\operatorname{Ind}_H^G W
=[G:H]\dim W.
\]

诱导表示最重要的性质是 Frobenius reciprocity：对任意 $G$-表示 $V$，有自然同构
\begin{equation}
\boxed{
\operatorname{Hom}_G(\operatorname{Ind}_H^G W,V)
\cong
\operatorname{Hom}_H(W,\operatorname{Res}_H^G V)
}.
\label{eq:frobenius-reciprocity}
\end{equation}
它说明“先把 $W$ 诱导到 $G$ 再寻找 $G$-intertwiner”与“先把 $V$ 限制到 $H$ 再寻找 $H$-intertwiner”是完全同一信息。


采用等价的张量模型
\[
\operatorname{Ind}_H^G W
=\CC[G]\otimes_{\CC[H]}W.
\]
若
\[
A:\CC[G]\otimes_{\CC[H]}W\to V
\]
是 $G$-intertwiner，定义
\[
\Phi(A):W\to V,
\qquad
\Phi(A)(w):=A(1\otimes w).
\]
对 $h\in H$，利用平衡张量积关系
\[
h\otimes w=1\otimes\sigma(h)w
\]
以及 $G$-equivariance，得到
\[
\begin{aligned}
\Phi(A)(\sigma(h)w)
&=A(1\otimes\sigma(h)w)\\
&=A(h\otimes w)\\
&=\rho(h)A(1\otimes w)\\
&=\rho(h)\Phi(A)(w),
\end{aligned}
\]
所以 $\Phi(A)$ 是 $H$-intertwiner。

反过来，若 $B:W\to\operatorname{Res}_H^G V$ 是 $H$-intertwiner，定义
\[
\Psi(B)(g\otimes w)
:=\rho(g)B(w).
\]
必须先验证它尊重平衡关系。对 $h\in H$，
\[
\begin{aligned}
\Psi(B)(gh\otimes w)
&=\rho(g)\rho(h)B(w)\\
&=\rho(g)B(\sigma(h)w)\\
&=\Psi(B)(g\otimes\sigma(h)w).
\end{aligned}
\]
故良定义。再有
\[
\Psi(B)(g_0g\otimes w)
=\rho(g_0)\Psi(B)(g\otimes w),
\]
所以它是 $G$-intertwiner。最后直接检查
\[
\Phi\circ\Psi=\IdMap,
\qquad
\Psi\circ\Phi=\IdMap,
\]
得到\eqnrefs{eq:frobenius-reciprocity}。


取不可约表示 $W=\tau$、$V=\rho$，Frobenius reciprocity 立刻给出重数公式
\[
\dim\operatorname{Hom}_G(\operatorname{Ind}_H^G\tau,\rho)
=
\dim\operatorname{Hom}_H(\tau,\operatorname{Res}_H^G\rho).
\]
因此“某个 $G$-不可约表示在诱导表示中出现几次”可以完全转化为“把它限制到 $H$ 后包含 $\tau$ 几次”。这在晶体小群、角动量耦合和对称破缺分支规则中都是基础结构。

接下来考察对称幂、外幂与置换群投影。

对单个表示 $V$，高阶张量 $V^{\otimes r}$ 同时携带 $G$ 的对角作用与置换群 $S_r$ 的因子置换作用：
\[
\rho^{\otimes r}(g)
(v_1\otimes\cdots\otimes v_r)
=
\rho(g)v_1\otimes\cdots\otimes\rho(g)v_r,
\]
\[
P_\pi(v_1\otimes\cdots\otimes v_r)
=
v_{\pi^{-1}(1)}\otimes\cdots\otimes v_{\pi^{-1}(r)}.
\]
两种作用彼此交换。于是可以先用 $S_r$ 的 idempotent 投影张量对称类型，再研究剩余 $G$-表示。

最基本的两个投影是
\begin{equation}
P_{\mathrm{sym}}
=
\frac1{r!}\sum_{\pi\in S_r}P_\pi,
\qquad
P_{\mathrm{alt}}
=
\frac1{r!}\sum_{\pi\in S_r}
\operatorname{sgn}(\pi)P_\pi.
\label{eq:symmetric-alternating-projectors}
\end{equation}
它们满足
\[
P_{\mathrm{sym}}^2=P_{\mathrm{sym}},
\qquad
P_{\mathrm{alt}}^2=P_{\mathrm{alt}},
\qquad
P_{\mathrm{sym}}P_{\mathrm{alt}}=0
\]
（当 $r>1$），其像分别是
\[
\operatorname{Sym}^rV,
\qquad
\Lambda^rV.
\]
因为 $P_\pi$ 与 $G$-作用交换，这两个子空间自动是 $G$-不变子表示。

当 $\dim V=n$ 时，
\[
\dim\operatorname{Sym}^rV
=\binom{n+r-1}{r},
\qquad
\dim\Lambda^rV
=\binom nr.
\]
特别地，$\Lambda^nV$ 是一维表示，其作用就是 determinant character：
\[
g\cdot(v_1\wedge\cdots\wedge v_n)
=
\det(\rho(g))
(v_1\wedge\cdots\wedge v_n).
\]
这说明 determinant 不是矩阵公式附带出来的量，而是最高外幂表示本身。

更一般的 Young symmetrizer 会对应 $S_r$ 的其他不可约表示，并把 $V^{\otimes r}$ 分解为不同张量对称类型。Schur--Weyl duality 进一步说明，在 $V=\CC^n$ 上，$GL(V)$ 与 $S_r$ 的交换作用彼此生成对方的 commutant；这为一般张量表示、Young diagram 与多体交换对称性提供统一来源。

\textbf{过渡：Lie 群、Lie 代数与 Casimir.}\enspace
\label{sec:lie-groups-algebras}

接下来考察从光滑群到 Lie 代数。

Lie 群 \(G\) 同时是光滑流形和群，并要求乘法与取逆映射光滑。单位元处切空间
\[
\mathfrak g=T_eG
\]
通过左平移可与所有左不变向量场一一对应。向量场的 Lie 括号因而在 \(T_eG\) 上诱导双线性运算
\[
[\cdot,\cdot]:\mathfrak g\times\mathfrak g\to\mathfrak g,
\]
满足反对称性和 Jacobi 恒等式
\begin{equation}
[X,[Y,Z]]+[Y,[Z,X]]+[Z,[X,Y]]=0.
\label{eq:jacobi-lie-algebra}
\end{equation}
由此得到 Lie 代数 \(\mathfrak g\)。

每个 \(X\in\mathfrak g\) 产生唯一的一参数子群
\[
\gamma_X(t+s)=\gamma_X(t)\gamma_X(s),
\qquad
\dot\gamma_X(0)=X.
\]
定义指数映射
\[
\exp:\mathfrak g\to G,
\qquad
\exp X=\gamma_X(1).
\]
矩阵 Lie 群中它就是矩阵指数。局部乘法不是简单的 \(\exp X\exp Y=\exp(X+Y)\)，非交换性由 Baker--Campbell--Hausdorff 展开控制：
\begin{equation}
\log(\ee^X\ee^Y)
=X+Y+\frac12[X,Y]
+\frac1{12}[X,[X,Y]]
+\frac1{12}[Y,[Y,X]]+\cdots.
\label{eq:bch-general}
\end{equation}



指数映射在单位元附近把 Lie algebra 的线性结构与群乘法联系起来。它的微分满足
\[
(\dd\exp)_0=\IdMap_{\mathfrak g},
\]
所以逆函数定理说明 $\exp$ 在 $0$ 的某个邻域内是到 $e$ 邻域的局部微分同胚。但它一般既不是全局一一，也不是群同态；失去交换性正由 BCH 的嵌套 commutator 测量。

一个特别重要的恒等式是
\begin{equation}
\boxed{
\operatorname{Ad}_{\exp X}
=\exp(\operatorname{ad}_X).
}
\label{eq:math-lie-ad-exp}
\end{equation}
这里右边是线性算子 $\operatorname{ad}_X:Y\mapsto[X,Y]$ 的普通指数。它把群级 conjugation 与 Lie algebra 级 commutator 精确连接起来。


固定 \(X,Y\in\mathfrak g\)，定义
\[
Y(t):=\operatorname{Ad}_{\exp(tX)}Y.
\]
这是 \(Y\) 沿 \(X\) 方向的共轭流。利用群同态性质
\[
\operatorname{Ad}_{gh}=\operatorname{Ad}_g\operatorname{Ad}_h
\]
（来自 \(C_{gh}=C_g\circ C_h\)，微分后即得）与前面得到的
\[
\left.\frac{\dd}{\dd s}\right|_{s=0}
\operatorname{Ad}_{\exp(sX)}Z
=[X,Z],
\]
有
\[
\begin{aligned}
\dot Y(t)
&=\lim_{h\to0}
\frac{\operatorname{Ad}_{\exp((t+h)X)}Y-
\operatorname{Ad}_{\exp(tX)}Y}{h}\\
&=\lim_{h\to0}
\frac{\operatorname{Ad}_{\exp(tX)}\operatorname{Ad}_{\exp(hX)}Y-
\operatorname{Ad}_{\exp(tX)}Y}{h}\\
&=\operatorname{Ad}_{\exp(tX)}
\lim_{h\to0}\frac{\operatorname{Ad}_{\exp(hX)}Y-Y}{h}\\
&=\operatorname{Ad}_{\exp(tX)}[X,Y]\\
&=[X,\operatorname{Ad}_{\exp(tX)}Y]\\
&=[X,Y(t)]
=\operatorname{ad}_X Y(t).
\end{aligned}
\]
倒数第二行用到 \(\operatorname{Ad}_{\exp(tX)}\) 是 Lie 代数同态，故与括号可交换：\(\operatorname{Ad}_{\exp(tX)}[X,Y]=[\operatorname{Ad}_{\exp(tX)}X,\operatorname{Ad}_{\exp(tX)}Y]=[X,Y(t)]\)，其中 \(\operatorname{Ad}_{\exp(tX)}X=X\) 因 \(\exp(tX)\) 与 \(X\) 对易。初值 \(Y(0)=\operatorname{Ad}_{e}Y=Y\)，所以这是常系数线性 ODE
\[
\dot Y=\operatorname{ad}_X Y,\qquad Y(0)=Y.
\]
按线性 ODE 唯一解定理，解为
\[
Y(t)=e^{t\operatorname{ad}_X}Y
=\sum_{n=0}^\infty\frac{t^n}{n!}\operatorname{ad}_X^n Y
=Y+t[X,Y]+\frac{t^2}{2}[X,[X,Y]]+\cdots.
\]
取 \(t=1\) 即得\eqnrefs{eq:math-lie-ad-exp}。

\emph{举例：\(\mathfrak{so}(3)\) 中的旋转。}取 \(X=\theta\hat n\cdot\vec J\)（\(\hat n\) 单位向量，\(J_i\) 角动量生成元满足 \([J_i,J_j]=\varepsilon_{ijk}J_k\)）。对任意 \(Y\in\mathfrak{so}(3)\)，
\[
\operatorname{Ad}_{\exp(\theta\hat n\cdot\vec J)}Y
=e^{\theta\hat n\cdot\operatorname{ad}_{\vec J}}Y
=Y+\theta[\hat n\cdot\vec J,Y]+\frac{\theta^2}{2}[\hat n\cdot\vec J,[\hat n\cdot\vec J,Y]]+\cdots.
\]
若 \(Y=\vec v\cdot\vec J\) 对应向量 \(\vec v\)，则 \([\hat n\cdot\vec J,\vec v\cdot\vec J]=(\hat n\times\vec v)\cdot\vec J\)，级数正好给出 Rodrigues 旋转公式
\[
\vec v\mapsto \vec v\cos\theta+(\hat n\times\vec v)\sin\theta+\hat n(\hat n\cdot\vec v)(1-\cos\theta).
\]
即 \(e^{\operatorname{ad}_X}\) 在 \(\mathfrak{so}(3)\) 上的作用就是绕 \(\hat n\) 转 \(\theta\) 角的旋转。

\emph{物理应用：Heisenberg 群与平移--动量对易。}对 Heisenberg 群，\(\operatorname{Ad}\) 平凡，\(\operatorname{ad}\) 仅记录对易子。在量子力学中，时间演化算子 \(U(t)=e^{-iHt/\hbar}\) 对可观测量 \(A\) 的 Heisenberg 演化
\[
A(t)=U(t)^{-1}AU(t)=e^{iHt/\hbar}Ae^{-iHt/\hbar}
\]
正是 \(\operatorname{Ad}_{\exp(iHt/\hbar)}A=e^{t\operatorname{ad}_{iH/\hbar}}A\)，其微分形式即 Heisenberg 方程
\[
\dot A=\frac{i}{\hbar}[H,A].
\]
此恒等式把 Schrödinger 图景与 Heisenberg 图景的等价性归结为 Lie 群--Lie 代数伴随关系的特例。

\emph{推广方向。}对一般 Banach Lie 群（如某些算子群），同一证明给出 \(\operatorname{Ad}_{\exp X}=e^{\operatorname{ad}_X}\)，是泛函分析中处理算子半群与扰动展开的代数基础。对量子群与 deformation quantization，此关系变形为 \(R\)-矩阵的 Yang--Baxter 方程，是量子可积系统的母结构。
\end{derivation}

BCH 的前几项也可由这一恒等式系统获得。令
\[
Z(t):=\log(e^{tX}e^{tY})
\]
并作形式展开
\[
Z(t)=tZ_1+t^2Z_2+t^3Z_3+O(t^4).
\]
逐阶比较 $e^{Z(t)}$ 与 $e^{tX}e^{tY}$ 的非交换乘积，就得到
\[
Z_1=X+Y,
\qquad
Z_2=\frac12[X,Y],
\]
以及三阶双重 commutator。于是 Lie bracket 可以理解为“群乘法偏离向量加法的最低非平凡项”。

接下来考察伴随表示与结构常数。

群在自身上通过共轭作用
\[
C_g(h)=ghg^{-1}
\]
产生伴随表示
\[
\operatorname{Ad}_g
=(\dd C_g)_e:\mathfrak g\to\mathfrak g.
\]
对 \(g=\exp(tX)\) 在 \(t=0\) 处求导得到 Lie 代数伴随表示
\[
\operatorname{ad}_X Y=[X,Y].
\]
在基 \(\{T_a\}\) 下定义结构常数
\[
[T_a,T_b]=f_{ab}{}^cT_c.
\]
Jacobi 恒等式等价于
\begin{equation}
f_{ab}{}^d f_{dc}{}^e
+f_{bc}{}^d f_{da}{}^e
+f_{ca}{}^d f_{db}{}^e=0.
\label{eq:structure-constant-jacobi}
\end{equation}
这不是某个矩阵群额外满足的巧合，而是所有 Lie 代数括号必须满足的相容条件。

Lie 代数表示是线性映射
\[
\dd\rho:\mathfrak g\to\operatorname{End}(V)
\]
满足
\[
[\dd\rho(X),\dd\rho(Y)]
=\dd\rho([X,Y]).
\]
若群表示 \(\rho:G\to GL(V)\) 光滑，则
\[
\dd\rho(X)
=\left.\frac{\dd}{\dd t}\rho(\exp tX)\right|_{t=0}.
\]
局部表示理论因此由生成元及其对易关系控制。

\begin{derivation}[由群共轭推出伴随 Lie 代数作用]
取矩阵 Lie 群以看清局部代数。对 \(X,Y\in\mathfrak g\)，由 \(\exp(tX)\in G\) 且共轭保持 Lie 代数，
\[
\operatorname{Ad}_{\exp(tX)}Y
=
(\dd C_{\exp(tX)})_eY
=
\exp(tX)Y\exp(-tX)
=
\ee^{tX}Y\ee^{-tX}.
\]
将矩阵指数写成级数 \(\ee^{tX}=\sum_{n=0}^\infty(tX)^n/n!\)，对乘积 \(\ee^{tX}Y\ee^{-tX}\) 用 Leibniz 乘积求导律，
\begin{align*}
\frac{\dd}{\dd t}
\bigl(\ee^{tX}Y\ee^{-tX}\bigr)
&=\left(\frac{\dd}{\dd t}\ee^{tX}\right)Y\ee^{-tX}
+\ee^{tX}Y\left(\frac{\dd}{\dd t}\ee^{-tX}\right)\\
&=X\ee^{tX}Y\ee^{-tX}
+\ee^{tX}Y(-X)\ee^{-tX}\\
&=X\ee^{tX}Y\ee^{-tX}
-\ee^{tX}Y\ee^{-tX}X.
\end{align*}
此处 \(\frac{\dd}{\dd t}\ee^{tX}=X\ee^{tX}\) 来自级数逐项求导：
\[
\frac{\dd}{\dd t}\sum_{n=0}^\infty\frac{t^nX^n}{n!}
=\sum_{n=1}^\infty\frac{t^{n-1}X^n}{(n-1)!}
=X\sum_{n=1}^\infty\frac{t^{n-1}X^{n-1}}{(n-1)!}
=X\ee^{tX},
\]
矩阵 \(X\) 与 \(\ee^{tX}\) 对易因 \(X\) 与自身各幂对易。类似地 \(\frac{\dd}{\dd t}\ee^{-tX}=-X\ee^{-tX}\)。在 \(t=0\) 处 \(\ee^{0}=I\)，故
\begin{align*}
\left.\frac{\dd}{\dd t}
\operatorname{Ad}_{\exp(tX)}Y
\right|_{t=0}
&=XIYI^{-1}-IYI^{-1}X\\
&=XY-YX\\
&=[X,Y].
\end{align*}
由 \(\operatorname{ad}_X\) 的定义 \(\operatorname{ad}_XY=[X,Y]\)，上式即
\[
\left.\frac{\dd}{\dd t}
\operatorname{Ad}_{\exp(tX)}
\right|_{t=0}
=\operatorname{ad}_X,
\]
因此伴随群表示 \(\operatorname{Ad}:G\to GL(\mathfrak g)\) 的微分正是 \(\operatorname{ad}:\mathfrak g\to\operatorname{End}(\mathfrak g)\)。同样的结论对一般 Lie 群由左不变向量场定义给出：取 \(Y\) 对应的左不变向量场 \(\widetilde Y\)，则 \(\operatorname{Ad}_gY\) 是 \((C_g)_*\widetilde Y\) 在单位元处的取值，对 \(g=\exp(tX)\) 求导得 Lie 括号 \([X,Y]\)，与矩阵实现无关。

\emph{举例：\(\mathfrak{gl}(n,\CC)\) 的伴随作用。}取 \(G=GL(n,\CC)\)，\(\mathfrak g=\mathfrak{gl}(n,\CC)=M_n(\CC)\)。对 \(X,Y\in M_n(\CC)\)，
\[
\operatorname{Ad}_{\exp(tX)}Y=e^{tX}Ye^{-tX},
\]
展开到二阶
\[
e^{tX}Ye^{-tX}=Y+t[X,Y]+\frac{t^2}{2}[X,[X,Y]]+O(t^3).
\]
若 \(X\) 为对角矩阵 \(\operatorname{diag}(\lambda_1,\ldots,\lambda_n)\)、\(Y=E_{ij}\)（矩阵单位），则 \([X,Y]=(\lambda_i-\lambda_j)E_{ij}\)，于是
\[
\operatorname{Ad}_{\exp(tX)}E_{ij}=e^{t(\lambda_i-\lambda_j)}E_{ij}.
\]
即伴随作用在矩阵单位上是对角的，本征值恰为根 \(\alpha_{ij}=\lambda_i-\lambda_j\)。这是 root decomposition 在 \(GL(n)\) 情形的直接体现。

\emph{物理应用：规范场的几何解释。}在 Yang--Mills 理论中，规范势 \(A_\mu\) 取值于 Lie 代数 \(\mathfrak g\)，规范变换 \(g(x)\) 通过伴随作用作用于势：
\[
A_\mu\mapsto gA_\mu g^{-1}-\frac{i}{g}(\partial_\mu g)g^{-1}.
\]
第一项正是 \(\operatorname{Ad}_g A_\mu\)，无穷小变换 \(g=\exp(tX)\) 给出 \(\delta A_\mu=t[X,A_\mu]=t\,\operatorname{ad}_X A_\mu\)。规范场强 \(F_{\mu\nu}\) 的协变性同样来自伴随作用。此推导把“规范变换如何作用于势”严格归结为 Lie 群--Lie 代数伴随关系。

\emph{推广方向。}对无穷维 Lie 群（如微分同胚群 \(\operatorname{Diff}(M)\)、规范群 \(\mathcal G=\operatorname{Map}(M,G)\)），伴随作用的微分仍是 Lie 括号，但需处理拓扑与定义域的细致问题。在 Kac--Moody 代数与顶点算子代数中，伴随关系扩展为带中心项的 projective 表示，是共形场论与弦论的代数基础。


接下来考察不变双线性型与 Casimir 元。

对半单 Lie 代数，Killing 型定义为
\[
B(X,Y)=\operatorname{tr}(\operatorname{ad}_X\operatorname{ad}_Y).
\]
它满足伴随不变性
\[
B([Z,X],Y)+B(X,[Z,Y])=0.
\]
若在某组基下不变非退化双线性型写作 \(\kappa_{ab}\)，逆矩阵为 \(\kappa^{ab}\)，则二次 Casimir 元可写成
\begin{equation}
C_2=\kappa^{ab}T_aT_b.
\label{eq:quadratic-casimir}
\end{equation}
利用不变性可验证 \([T_c,C_2]=0\)。因此在任意不可约复表示中，Schur 引理给出
\[
\dd\rho(C_2)=c_2(\rho)I.
\]
Casimir 本征值于是成为不可约表示的标签之一。



对紧致连通 Lie 群，还可以把表示论进一步压缩到一组彼此可交换的生成元。取 maximal torus $T\subset G$，其 Lie algebra 记为 $\mathfrak t$。由于 $\mathfrak t$ Abel，复化后的伴随算子
\[
\operatorname{ad}_H,
\qquad H\in\mathfrak t_\CC,
\]
可以同时对角化，于是半单情形得到 root decomposition
\begin{equation}
\boxed{
\mathfrak g_\CC
=\mathfrak t_\CC
\oplus
\bigoplus_{\alpha\in\Phi}\mathfrak g_\alpha,
}
\label{eq:math-lie-root-decomposition}
\end{equation}
其中
\begin{equation}
\mathfrak g_\alpha
:=\{X\in\mathfrak g_\CC:[H,X]=\alpha(H)X
\ \forall H\in\mathfrak t_\CC\}.
\label{eq:math-lie-root-space}
\end{equation}
非零线性泛函 $\alpha\in\mathfrak t_\CC^*$ 称为 root。Jacobi 恒等式立即给出
\begin{equation}
[\mathfrak g_\alpha,\mathfrak g_\beta]
\subseteq\mathfrak g_{\alpha+\beta},
\label{eq:math-lie-root-bracket}
\end{equation}
其中若 $\alpha+\beta$ 不是 root 且不为零，则右边理解为 $0$。

同样，对一个有限维复表示 $V$，可按 $\mathfrak t$ 的共同本征值分解为 weight spaces
\begin{equation}
V=\bigoplus_{\lambda}V_\lambda,
\qquad
V_\lambda
:=\{v:\dd\rho(H)v=\lambda(H)v
\ \forall H\in\mathfrak t_\CC\}.
\label{eq:math-lie-weight-decomposition}
\end{equation}
若 $X_\alpha\in\mathfrak g_\alpha$，则
\begin{equation}
\boxed{
\dd\rho(X_\alpha)V_\lambda
\subseteq V_{\lambda+\alpha}.
}
\label{eq:math-lie-root-shifts-weight}
\end{equation}
所以 root operators 是 weight diagram 上的“升降方向”。


取 \(v\in V_\lambda\)，由 weight space 定义对任意 \(H\in\mathfrak t_\CC\) 有
\[
\dd\rho(H)v=\lambda(H)v.
\]
因 \(X_\alpha\in\mathfrak g_\alpha\)，由 root space 定义 \([H,X_\alpha]=\alpha(H)X_\alpha\)。又因 \(\dd\rho\) 是 Lie 代数同态，它保持括号：
\[
[\dd\rho(H),\dd\rho(X_\alpha)]
=\dd\rho(H)\dd\rho(X_\alpha)-\dd\rho(X_\alpha)\dd\rho(H)
=\dd\rho([H,X_\alpha])
=\alpha(H)\dd\rho(X_\alpha).
\]
把换位子展开并移项，得
\[
\dd\rho(H)\dd\rho(X_\alpha)v
=\dd\rho(X_\alpha)\dd\rho(H)v
+\alpha(H)\dd\rho(X_\alpha)v.
\]
把 \(\dd\rho(H)v=\lambda(H)v\) 代入右端第一项，
\begin{align*}
\dd\rho(H)\dd\rho(X_\alpha)v
&=\dd\rho(X_\alpha)\{\lambda(H)v\}
 +\alpha(H)\dd\rho(X_\alpha)v\\
&=\lambda(H)\dd\rho(X_\alpha)v
 +\alpha(H)\dd\rho(X_\alpha)v\\
&=(\lambda(H)+\alpha(H))\dd\rho(X_\alpha)v.
\end{align*}
上式对任意 \(H\in\mathfrak t_\CC\) 成立。由 weight space \(V_{\lambda+\alpha}\) 的定义，这意味着只要 \(\dd\rho(X_\alpha)v\neq0\)，它就是权 \(\lambda+\alpha\) 的权向量，属于 \(V_{\lambda+\alpha}\)。因此
\[
\dd\rho(X_\alpha)V_\lambda\subseteq V_{\lambda+\alpha},
\]
即\eqnrefs{eq:math-lie-root-shifts-weight}。

\emph{举例：\(\mathfrak{su}(2)\) 的升降算子。}取 \(\mathfrak{su}(2)_\CC=\mathfrak{sl}(2,\CC)\)，Cartan 子代数由 \(H=\operatorname{diag}(1,-1)\) 张成。root 为 \(\alpha(H)=2\) 与 \(-\alpha(H)=-2\)，相应 root operators
\[
E=\begin{pmatrix}0&1\\0&0\end{pmatrix},\qquad
F=\begin{pmatrix}0&0\\1&0\end{pmatrix}.
\]
在自旋 \(j\) 表示中，权为 \(\lambda_m(H)=2m\)（\(m=-j,-j+1,\ldots,j\)）。公式给出
\[
E V_{\lambda_m}\subseteq V_{\lambda_{m+1}},\qquad
F V_{\lambda_m}\subseteq V_{\lambda_{m-1}},
\]
即 \(E\) 为升算子、\(F\) 为降算子。具体地 \(E|m\rangle\propto|m+1\rangle\)、\(F|m\rangle\propto|m-1\rangle\)，比例系数由 \(\mathfrak{sl}(2)\) 表示论确定为 \(\sqrt{j(j+1)-m(m\pm1)}\)。这正是量子力学角动量阶梯算子的代数来源。

\emph{举例：\(\mathfrak{su}(3)\) 与夸克模型。}\(\mathfrak{su}(3)_\CC=\mathfrak{sl}(3,\CC)\) 有两个简单 root \(\alpha_1,\alpha_2\)，Cartan 由 \(H_1=\operatorname{diag}(1,-1,0)\)、\(H_2=\operatorname{diag}(0,1,-1)\) 张成。基本表示（三维）的权为
\[
\lambda_1=(\tfrac23,-\tfrac13),\quad
\lambda_2=(-\tfrac13,\tfrac23),\quad
\lambda_3=(-\tfrac13,-\tfrac13),
\]
对应上、下、奇异三种夸克味道。root operators \(E_{\alpha_1},E_{\alpha_2},E_{\alpha_1+\alpha_2}\) 把一种夸克“变为”另一种，正是 Gell--Mann \(\lambda\) 矩阵的非对角部分。八重态与十重态的权图即由这些 root 平移操作从最高权出发逐步生成。

\emph{物理应用：角动量代数的建模链条。}从 Lie 代数 \(\mathfrak{su}(2)\) 出发：(1) Cartan 元 \(J_z\) 给出权（磁量子数 \(m\)）；(2) root operators \(J_\pm=J_x\pm iJ_y\) 为升降算子；(3) 最高权 \(j\) 标记不可约表示；(4) 物理可观测量 \(J_z,J^2\) 对应到 Cartan 与 Casimir；(5) 对称性约束（转动不变性）要求能级按 \(SU(2)\) 不可约表示分类，给出 \(2j+1\) 重简并。整个链条从 root--weight 代数结构通向原子光谱与自旋系统的简并模式。

\emph{推广方向。}对 Kac--Moody 代数（仿 Lie 代数），同一证明给出 root operator 在 weight 上的平移，但 root 系统可含虚根，权图变为无限。对量子群 \(U_q(\mathfrak g)\)，平移关系变形为 \(q\)-commutator，给出量子 Clebsch--Gordan 与 crystal base 理论。在几何表示论中，Hecke 代数与 Springer 理论把 root--weight 组合解释为旗流形上的几何操作，是 Langlands 纲领的代数支柱。
\end{derivation}

选择一组 positive roots 后，每个有限维不可约表示都有 highest weight；对紧致半单群，highest-weight theory 把不可约表示的分类转化为允许的 integral dominant weights。这里“integral”不仅取决于 Lie algebra，还取决于群的全局拓扑：具有相同 Lie algebra 的 $SU(2)$ 与 $SO(3)$ 允许的群表示并不完全相同。Weyl group
\[
W=N_G(T)/T
\]
则作用在 $\mathfrak t^*$ 上并置换 roots/weights，编码 maximal torus 选择的剩余离散对称性。

这条结构链把有限群 character theory 与连续群表示论对应起来：有限群中用 conjugacy classes 和 characters 把群代数对角化；紧致 Lie 群中则先对 maximal torus 同时对角化，再用 roots 记录非交换方向如何连接不同 weights。

接下来考察一般理论后的例子：\(SU(2)\) 与 \(SO(3)\)。

\(SU(2)\) 的 Lie 代数可取 Hermitian 生成元
\[
T_i=\frac12\sigma_i,
\qquad
[T_i,T_j]=\ii\varepsilon_{ijk}T_k,
\]
其中 \(\sigma_i\) 为 Pauli 矩阵。抽象地，这只是三维实 Lie 代数 \(\mathfrak{su}(2)\) 的一个二维复表示。其二次 Casimir
\[
C_2=T_1^2+T_2^2+T_3^2
\]
在自旋 \(j\) 不可约表示中取值
\[
C_2=j(j+1)I,
\qquad
\dim V_j=2j+1.
\]

群级别存在满同态
\[
SU(2)\longrightarrow SO(3)
\]
且核为 \(\{\pm I\}\)，所以
\[
SO(3)\cong SU(2)/\{\pm I\}.
\]
这说明半整数 \(j\) 的表示是 \(SU(2)\) 的诚实表示，却只给出 \(SO(3)\) 的投影表示。下一章从 Clifford 代数的 universal property 出发，会把这种二重覆盖重新解释为一般的 \(\operatorname{Spin}(n)\to SO(n)\) 结构，而不是把 Pauli 矩阵当作起点。

接下来考察从 Lie 代数看 $SU(3)$ 的八个生成方向。

上面的八个 Gell--Mann 矩阵不仅是一张矩阵表。它们恰好把 ``为什么是八个自由方向''、``有限变换为什么由指数产生'' 以及 ``三分量态怎样在无穷小变换下混合'' 连在一起。

先从群的约束计数。一般复 \(3\times3\) 矩阵有 \(18\) 个实参数。幺正条件
\[
U^\dagger U=I
\]
给出一个 Hermitian 矩阵方程，相当于 \(9\) 个独立实约束，所以 \(U(3)\) 有 \(18-9=9\) 个实参数。再施加
\[
\det U=1
\]
去掉一个连续相位自由度，便得到
\[
\dim_{\mathbb R}SU(3)=8.
\]
这与八个 \(T_a=\lambda_a/2\) 完全一致。

为了从有限群元素反推出 Lie 代数，取单位元附近的一条光滑曲线
\[
U(\varepsilon)=I-\ii\varepsilon H+O(\varepsilon^2).
\]
由幺正条件
\[
U^\dagger U
=I+\ii\varepsilon(H^\dagger-H)+O(\varepsilon^2)=I
\]
得到
\[
H^\dagger=H.
\]
另一方面
\[
1=\det U(\varepsilon)
=\exp\!\bigl(\operatorname{tr}\ln U\bigr)
=1-\ii\varepsilon\operatorname{tr}H+O(\varepsilon^2),
\]
故
\[
\operatorname{tr}H=0.
\]
因此物理约定下的无穷小生成元就是全部 Hermitian、无迹的 \(3\times3\) 矩阵，而 \(T_a\) 正好构成这个八维实向量空间的一组正交基。

这一点还给出一个方便的系数抽取公式。若
\[
H=h_aT_a,
\]
利用
\[
\operatorname{tr}(T_aT_b)=\frac12\delta_{ab},
\]
可得
\[
\boxed{
h_a=2\operatorname{tr}(HT_a)
}.
\]
所以任意无迹 Hermitian 三阶矩阵都可以唯一地按 Gell--Mann 基分解。

结构常数也不是抽象地“给出”的。由
\[
[T_a,T_b]=\ii f_{abc}T_c
\]
两边乘 \(T_d\) 后取迹，得到
\[
\operatorname{tr}([T_a,T_b]T_d)
=\frac{\ii}{2}f_{abd},
\]
因而
\[
\boxed{
f_{abd}=-2\ii\operatorname{tr}([T_a,T_b]T_d)
}.
\]
例如
\[
\lambda_1\lambda_2
=\begin{pmatrix}\ii&0&0\\0&-\ii&0\\0&0&0\end{pmatrix}
=\ii\lambda_3,
\]
而
\[
\lambda_2\lambda_1=-\ii\lambda_3,
\]
所以
\[
[\lambda_1,\lambda_2]=2\ii\lambda_3,
\qquad
[T_1,T_2]=\ii T_3,
\]
即 \(f_{123}=1\)。这与 Pauli 矩阵子块中的 \(SU(2)\) 结构完全一致。

\(T_3\) 与 \(T_8\) 都是对角矩阵，而且
\[
[T_3,T_8]=0.
\]
它们张成 \(\mathfrak{su}(3)\) 的一个 Cartan 子代数。对三维基本表示的三个标准基矢
\[
e_1=\begin{pmatrix}1\\0\\0\end{pmatrix},\qquad
 e_2=\begin{pmatrix}0\\1\\0\end{pmatrix},\qquad
 e_3=\begin{pmatrix}0\\0\\1\end{pmatrix},
\]
有
\[
T_3e_1=\frac12e_1,
\qquad
T_8e_1=\frac1{2\sqrt3}e_1,
\]
\[
T_3e_2=-\frac12e_2,
\qquad
T_8e_2=\frac1{2\sqrt3}e_2,
\]
\[
T_3e_3=0,
\qquad
T_8e_3=-\frac1{\sqrt3}e_3.
\]
因此三个基矢可以按一对同时本征值标记为
\[
\left(\frac12,\frac1{2\sqrt3}\right),\qquad
\left(-\frac12,\frac1{2\sqrt3}\right),\qquad
\left(0,-\frac1{\sqrt3}\right).
\]
这就是三维基本表示的三个权。其余六个非对角生成元恰好在这些权之间作跃迁。例如
\[
(T_1+iT_2)e_2=e_1,
\]
而
\[
(T_1-iT_2)e_1=e_2.
\]
从这个角度看，Gell--Mann 矩阵不是八个彼此无关的对象，而是“两条对角测量方向 + 三对升降方向”的完整 Lie 代数结构。

这也解释了为什么仅仅说“\(SU(3)\) 有八个生成元”还不够。真正有用的是知道生成元怎样闭合、哪些能同时对角化、哪些负责不同权态之间的跃迁，以及有限群元素怎样由
\[
U=\exp(-\ii\theta_aT_a)
\]
从这些无穷小方向积分出来。

